\section{讨论}
\label{sec:discussion}

理论结果是在明确观测模型下区分推断目标。报告固定已知、暴露规模外生且计数误差条件独立时，当期值风险为 $A/m$，已观测时期实现均值的风险为 $A/(Tm)$。估计平稳过程均值，还需计入该实现均值的方差。在固定时期内增加病例不会改变这一贡献，而随时期数增加时它如何变化取决于时间协方差。报告波动则会改变当期值与实现均值的风险。这些恒等式为解释病例量提供参照，其假设决定参照的适用范围。

这一方差结构在种群统计随机性、环境随机性与整群抽样中已有先例\citep{engen1998demographic, cochran1977sampling}。其监测含义取决于哪些量共同变化，以及估计意在描述什么。独立后代与独立报告产生随暴露规模下降的抽样误差；共同传播条件通过所抽取的时期影响过程均值估计；共同报告波动影响对这些时期的测量。把所有持续存在的分量都称为环境底线，会掩盖这一区别，并可能把目标本身的变异误当作测量误差。

条件信息计算与流行曲线信息量的分析相互补充\citep{parag2022quantifying, parag2022fundamental}。它确立了分母和报告分布已知时，尺度化计数在指定无偏估计量类中的有效性，并未确立再生数推断的普适下界。追加观测、独立报告测量与先验结构都会改变可用信息。以条件偏差换取较低方差的方法同样位于该估计量类之外。因此，平衡规模应始终附属于其估计量、目标与抽样实验，而不能视为病原体固有的属性。

另一个相关区别涉及估计窗口。在 Cori 框架中，扩大窗口可以累积感染压力与发病数，同时要求再生数在更长区间内局部恒定\citep{cori2013new, gostic2020practical}。这在时间分辨率与精度之间进行权衡，传播发生变化时可能引入模型失配。它不等同于抽取更多独立传播时期，也不会自动保持估计对象不变。表~\ref{tab:targets} 的重复时期平均使用另一种假设：每个时期各有其传播值，并有对该值条件无偏的测量。把这些公式用于重叠窗口时，需明确处理窗口重叠与时间相依性。

实证分析考察的是表观周统计量的相对偏差平方，而非相对于已观测 $R_t$ 的均方误差。第~\ref{sec:res-empirical}~节的仿射与幂律比较，涉及观测计数范围内得到支持的曲线形态。仅有正的拟合截距不能确立非零渐近极限，预测表现接近也不能识别生成任一曲线的机制。仅用训练样本确定中心与评分权重，可避免预测比较中的信息泄漏。前瞻曲线区分分析仍以其生成机制为条件，回答的问题不同于识别生物学传播方差。

受控稀释对已记录计数施加已知抽样操作。随机移除前的原始观测轨迹得到保留，但非线性变换、连续性校正与随机分母都可能改变两个拟合系数。基于导数的参照是对整个计算流程的近似，并非统计量的精确不变量。因此，解读第~\ref{sec:res-thinning}~节时，必须同时考虑各保留率下的响应、不确定性与已知计数模型检查。与参照响应一致，可以支持该局部抽样描述，却不能验证住院数据满足完整的 $A/m$ 恒等式，也不能识别未稀释斜率的来源。未拒绝截距变化不构成等效性的证据。

州截距拟合与替代中心方案回答不同的问题。州截距允许各州基准损失不同，同时估计损失与州内计数变化的关联；它不会使这种变化成为外生变量，也不会消除共同周冲击。重抽样摘要必须对应目标州分布，并计入抽中次数；未抽中的州没有已估计截距。按流行季、趋势或全国周重新中心化，改变了被建模的损失。相应截距变化是敏感性指标，并非正交的生物学方差分量。全国周均值的任何方差比，针对的是方法中定义的有符号偏差，而不会自动对应平方损失响应。

空间与时间相依性也约束不确定性。同周地区受到共同影响，同一地区的重复观测具有共同历史。交叉重抽样在其自身假设下处理这两类相依性\citep{owen2007pigeonhole, cameron2011robust}，限定于季内的时间块敏感性分析则考察剩余序列相依。无论增加同一流行季内的数据行，还是增加住院数，都不会提供独立流行季。类似地，如式~(\ref{eq:pooling}) 所示，空间聚合可以通过局部流行之间的协方差改变传播本身的变异。因此，横截面计数关联可能同时包含抽样精度与被平均人群差异的影响。

模拟反演以传播和观测机制为条件（第~\ref{sec:meth-inversion}~节与第~\ref{sec:res-sim}~节）。每日更新模拟器中指定的周再生数变异，是按日历周索引的模型参数，并非直接观测到的逐代真值；对病例数乘以时间换算因子，也不能换算其方差或平衡规模。独立观测离散、共同报告冲击与传播持续性必须分别规定。由同一住院序列计算的统计量彼此一致，是有用的一致性检查，但共同数据与共同校准假设限制了它作为证据的独立性。

代际间隔进一步限定这一映射。序列间隔、内禀代际间隔与家庭内实现代际间隔衡量不同的量\citep{cowling2009estimation, vink2014serial, chan2025estimating}。采用假定变异系数构造的敏感性分布，必须与来源研究报告的离散程度区分；间隔均值的置信限也不定义整个代际间隔分布的后验。条件校准区间与指定情景之间的变化描述不同的不确定性，情景范围不是联合置信区间。传播不确定性需要有依据的分布与明确的联合模型，不能靠给敏感性端点重新命名来实现。

最后，机构参与度只衡量检出过程的一方面。参与报告的医院内部，检测、就医、编码与入院实践也可能变化。缺乏识别政策效应的设计时，报告制度间比较仍属描述性分析。对这些过程的独立测量、更细空间观测以及更多流行季，可以直接应对主要识别限制。附录~\ref{app:extensions} 的符号化设计扩展示范所选目标如何进入资源分配，同时受成本、可行时期数和相依性约束。实际决策要求在一致尺度上测量这些量。单凭经验摘要，无法给出普适病例门槛或唯一识别的传播方差。
